%% ── 2-Р БҮЛЭГ: ОНОЛЫН ХЭСЭГ ─────────────────────────────────
\chapter{Судалгааны Ажлын Онолын Хэсэг, Өнөөгийн Түвшин}
\label{ch:literature}

\section{Тэмдэглэл хөтлөлтийн түүхэн хөгжил}

Тэмдэглэл хөтлөх үйл ажиллагаа нь хүний бичгийн соёл үүссэн цагаас эхлэлтэй олон
мянган жилийн түүхтэй оюуны үйл ажиллагаа юм.

\subsection{Дижитал тэмдэглэлийн эхлэл (1990--2000)}

Компьютер, интернэт түгээмэл болсноор тэмдэглэл дараах хэлбэрт шилжсэн:
\begin{itemize}
  \item Microsoft OneNote (2003) -- хэвлэлийн байгууллагуудад түгээмэл
  \item Evernote (2008) -- хоорондоо синхрончлогдох анхны систем
  \item Google Keep (2013) -- хурдан тэмдэглэлд зориулагдсан, хөнгөн
  \item Notion (2016) -- бүрэн workspace платформ
\end{itemize}

\subsection{Хиймэл оюун ухаан суурилсан тэмдэглэлийн үе (2020--өнөөдөр)}

Сүүлийн жилүүдэд AI технологи нэмэгдсэнээр тэмдэглэлийн системүүд дараах боломжтой болсон:
\begin{itemize}
  \item Текстийг автоматаар эмхэтгэх, хураангуй гаргах
  \item Тэмдэглэлийн агуулгаар зураг, sticker үүсгэх
  \item Sentiment analysis ашиглан хэрэглэгчийн сэтгэл хөдлөл тодорхойлох
  \item Автомат мэдэгдэл, сануулга илгээх
\end{itemize}

\section{Дижитал тэмдэглэлийн системийн онолын үндэс}

\subsection{REST API архитектур}

REST (Representational State Transfer) нь HTTP протоколд суурилсан архитектурын загвар
бөгөөд client-server харилцааг тогтмолжуулдаг. REST API-ийн үндсэн зарчмууд:
\begin{itemize}
  \item \textbf{Stateless} -- сервер нь client-ийн төлөвийг хадгалдаггүй
  \item \textbf{Uniform Interface} -- нэгдсэн endpoint бүтэц (GET, POST, PUT, DELETE)
  \item \textbf{Client-Server} -- frontend болон backend тусдаа хөгжүүлэгддэг
  \item \textbf{Cacheable} -- хариуг cache хийж гүйцэтгэлийг сайжруулах
\end{itemize}

\subsection{JWT (JSON Web Token) баталгаажуулалт}

JWT нь хэрэглэгчийн баталгаажуулалтад ашиглагддаг стандарт токен форматыг хэлнэ~\cite{jwt2015}.
Бүтэц нь 3 хэсгээс тогтоно: \textbf{Header.Payload.Signature}

\begin{itemize}
  \item \textbf{Header} -- алгоритмын мэдээлэл (HS256)
  \item \textbf{Payload} -- хэрэглэгчийн мэдээлэл (userId, expiresIn)
  \item \textbf{Signature} -- нууц түлхүүрээр гарын үсэг зурсан хэсэг
\end{itemize}

Манай системд JWT токеныг \texttt{flutter\_secure\_storage}-д аюулгүй хадгалж,
дараагийн API дуудалтад \texttt{Authorization: Bearer <token>} header-т дамжуулна.

\subsection{MongoDB NoSQL өгөгдлийн сан}

MongoDB нь document-oriented NoSQL өгөгдлийн сан бөгөөд JSON-тай ижил BSON форматаар
өгөгдөл хадгалдаг~\cite{mongodb2019}. Давуу талууд:
\begin{itemize}
  \item Schema-flexible -- өгөгдлийн бүтцийг уян хатнаар өөрчлөх боломжтой
  \item Horizontal scaling -- өгөгдлийг олон серверт тарааж хадгалах
  \item Aggregation pipeline -- хүчирхэг шүүлт, нэгтгэлтийн боломж
  \item Atlas -- үүлэн хостолт, автомат backup
\end{itemize}

\subsection{Flutter мобайл хөгжүүлэлтийн платформ}

Flutter нь Google-ийн гаргасан нээлттэй эхийн UI toolkit бөгөөд нэг код бичиж
Android, iOS, Web, Desktop платформуудад ажиллах апп хөгжүүлэх боломжийг олгодог~\cite{windmill2019}.

\begin{itemize}
  \item \textbf{Dart хэл} -- OOP, strongly-typed, AOT compiled
  \item \textbf{Widget tree} -- UI-г widget-үүдийн мод хэлбэрээр байгуулдаг
  \item \textbf{Hot Reload} -- кодын өөрчлөлтийг тэр даруй харах
  \item \textbf{Material Design} -- Google-ийн дизайн стандартыг дэмжих
\end{itemize}

\section{Mood Tracking технологийн үндэс}

Mood tracking буюу сэтгэлийн байдлыг хянах нь хэрэглэгчийн сэтгэл хөдлөлийг
тодорхойлж, тэдэнд тохирсон контент, сануулга гаргах технологи юм.

\begin{longtable}{|p{3cm}|p{5cm}|p{5cm}|}
\caption{Mood tracking-ийн төрлүүд}\\
\hline
\textbf{Mood} & \textbf{Дүрлэл} & \textbf{Тайлбар} \\
\hline
Happy & Emoji: :) -- аз жаргалтай & Баяр хөөртэй, эерэг сэтгэлийн байдал \\
\hline
Sad & Emoji: :( -- гунигтай & Гунигтай, уналтын сэтгэлийн байдал \\
\hline
Tired & Emoji: -\_- -- ядарсан & Ядарсан, унтмаар байгаа байдал \\
\hline
Angry & Emoji: $>$:( -- уурласан & Уурласан, хурцадмал байдал \\
\hline
Love & Emoji: $<$3 -- дуртай & Дуртай, халуун дотно байдал \\
\hline
\end{longtable}

\section{Төстэй бүтээгдэхүүнүүдтэй харьцуулалт}

\begin{longtable}{|p{3.5cm}|p{2.5cm}|p{2.5cm}|p{2.5cm}|p{2.5cm}|}
\caption{Төстэй бүтээгдэхүүнүүдтэй харьцуулалт}\\
\hline
\textbf{Үзүүлэлт} & \textbf{Манай апп} & \textbf{Notion} & \textbf{Evernote} & \textbf{Google Keep} \\
\hline
CRUD үйлдэл & \checkmark & \checkmark & \checkmark & \checkmark \\
\hline
Mood tracking & \checkmark & -- & -- & -- \\
\hline
JWT auth & \checkmark & \checkmark & \checkmark & \checkmark \\
\hline
Монгол UI & \checkmark & -- & -- & -- \\
\hline
REST API & \checkmark & -- & \checkmark & -- \\
\hline
Нээлттэй эх & \checkmark & -- & -- & -- \\
\hline
\end{longtable}

\section{Express.js фреймворкийн онолын үндэс}

Express.js нь Node.js-д суурилсан хурдан, хялбар веб framework юм~\cite{nodejs2020}.
Middleware болон routing механизмаараа RESTful API барихад өргөн хэрэглэгддэг.

\subsection{Middleware дарааллын үзэл баримтлал}

Express.js-д middleware гэдэг нь HTTP хүсэлтийн дунд гүйцэтгэгддэг функц юм.
Дараалал дараах байдлаар явагдана:

\begin{enumerate}
  \item Клиентийн хүсэлт ирнэ (HTTP Request)
  \item \texttt{cors()} -- Cross-Origin аюулгүй байдлыг шалгана
  \item \texttt{express.json()} -- JSON body-г parse хийнэ
  \item \texttt{authMiddleware} -- JWT токеныг шалгана (хамгаалагдсан route-уудад)
  \item Route handler -- бизнесийн логик гүйцэтгэнэ
  \item HTTP хариу буцаана
\end{enumerate}

\subsection{Mongoose ODM (Object Document Mapper)}

Mongoose нь MongoDB-тэй ажиллахад schema-based шийдэл санал болгодог Node.js
library юм. Schema тодорхойлж, model үүсгэж, CRUD үйлдэл хийхэд ашиглана.

\begin{longtable}{|p{4cm}|p{9cm}|}
\caption{Mongoose ODM-ийн үндсэн аргуудын тайлбар}\\
\hline
\textbf{Mongoose арга} & \textbf{Тайлбар} \\
\hline
\texttt{Model.save()} & Шинэ document хадгалах \\
\hline
\texttt{Model.find()} & Нөхцөлийн дагуу хайх \\
\hline
\texttt{Model.findOneAndUpdate()} & Нэг document олж шинэчлэх \\
\hline
\texttt{Model.findOneAndDelete()} & Нэг document олж устгах \\
\hline
\texttt{Model.findOne()} & Нөхцөлд тохирсон эхний doc олох \\
\hline
\end{longtable}

\section{flutter\_secure\_storage-ийн онолын үндэс}

\texttt{flutter\_secure\_storage} нь Flutter-т зориулсан аюулгүй хадгалалтын package юм.
iOS-д Keychain, Android-д EncryptedSharedPreferences ашигладаг.

\begin{longtable}{|p{2.8cm}|p{5.8cm}|p{5.4cm}|}
\caption{flutter\_secure\_storage платформын харьцуулалт}\\
\hline
\textbf{Платформ} & \textbf{Технологи} & \textbf{Аюулгүй байдлын түвшин} \\
\hline
Android & EncryptedShared\-Preferences & AES 256-bit шифрлэлт \\
\hline
iOS & Keychain Services & Hardware-level аюулгүй байдал \\
\hline
\end{longtable}

Хадгалалтын үндсэн арга:
\begin{itemize}
  \item \texttt{write(key, value)} -- утга хадгалах
  \item \texttt{read(key)} -- утга унших
  \item \texttt{delete(key)} -- утга устгах
  \item \texttt{deleteAll()} -- бүгдийг устгах (logout)
\end{itemize}

\section{HTTP Client Flutter дахь ашиглалт}

Flutter-д \texttt{http} package ашиглан REST API-тай харьцана.
Үндсэн HTTP арга:

\begin{longtable}{|p{3cm}|p{4cm}|p{6cm}|}
\caption{Flutter HTTP Client аргуудын тайлбар}\\
\hline
\textbf{HTTP арга} & \textbf{Flutter код} & \textbf{Зориулалт} \\
\hline
GET & \texttt{http.get(uri, headers)} & Өгөгдөл авах \\
\hline
POST & \texttt{http.post(uri, body)} & Шинэ өгөгдөл үүсгэх \\
\hline
PUT & \texttt{http.put(uri, body)} & Өгөгдөл шинэчлэх \\
\hline
DELETE & \texttt{http.delete(uri)} & Өгөгдөл устгах \\
\hline
\end{longtable}

\section{CORS (Cross-Origin Resource Sharing)}

Манай систем frontend болон backend өөр портод ажилладаг тул CORS тохиргоо зайлшгүй
шаардлагатай. Backend-д \texttt{cors()} middleware-ийг \texttt{server.js}-д ашиглана.

CORS-ийн тохиргоогүй бол browser нь cross-origin хүсэлтийг хориглоно.
\texttt{npm install cors} командаар суулгаж, \texttt{app.use(cors())} дуудна.

\section{RESTful API дизайны зарчим}

Манай API-ийн endpoint-үүд Richardson Maturity Model-ийн 2-р түвшинд байна
(HTTP арга + Resource-based URL):

\begin{longtable}{|p{4cm}|p{3cm}|p{6cm}|}
\caption{RESTful API endpoint-үүдийн жагсаалт}\\
\hline
\textbf{Resource} & \textbf{URL} & \textbf{HTTP методууд} \\
\hline
Хэрэглэгч & /api/auth & POST (login, signup) \\
\hline
Тэмдэглэл & /api/notes & GET, POST \\
\hline
Тэмдэглэл (id-р) & /api/notes/:id & PUT, DELETE \\
\hline
\end{longtable}
